% to compile: Shift+Alt+F1
\documentclass[a4paper,11pt]{article}%,twocolumn
%% packages

\usepackage{blindtext} % needed for creating dummy text passages
%\usepackage{ngerman} % needed for German default language
\usepackage{amsmath} % needed for command eqref
\usepackage{amssymb} % needed for math fonts
\usepackage[colorlinks=true,breaklinks]{hyperref} % needed for creating hyperlinks in the document, the option colorlinks=true gets rid of the awful boxes, breaklinks breaks lonkg links (list of figures), and ngerman sets everything for german as default hyperlinks language
\usepackage[hyphenbreaks]{breakurl} % ben�tigt f�r das Brechen von URLs in Literaturreferenzen, hyphenbreaks auch bei links, die �ber eine Seite gehen (mit hyphenation).
\usepackage{xcolor}
\definecolor{c1}{rgb}{0,0,1} % blue
\definecolor{c2}{rgb}{0,0.3,0.9} % light blue
\definecolor{c3}{rgb}{0.3,0,0.9} % red blue
\hypersetup{
    linkcolor={c1}, % internal links
    citecolor={c2}, % citations
    urlcolor={c3} % external links/urls
}
%\usepackage{cite} % needed for cite
\usepackage[square,authoryear]{natbib} % needed for cite and abbrvnat bibliography style
\usepackage[nottoc]{tocbibind} % needed for displaying bibliography and other in the table of contents
\usepackage{graphicx} % needed for \includegraphics 
\usepackage{longtable} % needed for long tables over pages
\usepackage{bigstrut} % needed for the command \bigstrut
\usepackage{enumerate} % needed for some options in enumerate
%\usepackage{todonotes} % needed for todos
\usepackage{makeidx} % needed for creating an index
\makeindex
\usepackage{gensymb}
%\usepackage{url}
\usepackage{xurl}
\usepackage{psfrag}
\usepackage{multirow}
\usepackage{subfigure}

\usepackage{algpseudocode}
\usepackage{float}
%\usepackage{minted}
\usepackage{tcolorbox}
\usepackage{menukeys}
\usepackage[
width=.8\textwidth,
justification=centering]{caption}
\usepackage{xcolor}
\usepackage{pstricks}
\usepackage{acronym}
%% page settings

\usepackage[top=20mm, bottom=20mm,left=20mm,right=20mm]{geometry} % needed for page border settings
\parindent=0mm % for space of first line of new text block
\sloppy % for writing with hyphenless justification (tries to)
\hyphenation{} % use hyphenation of tolerance parametershttp://www.jr-x.de/publikationen/latex/tipps/zeilenumbruch.html
\hyphenpenalty=10000
\exhyphenpenalty=10000
\usepackage{fancyhdr} % needed for head and foot options
%% my macros

%% Text fomats
\newcommand{\tbi}[1]{\textbf{\textit{#1}}}

% IEEE naming convention
%\renewcommand{\figurename}{Fig.}

% custom counter for shell scripts
\newcounter{shellcounter}
% \newcounter{myexample}% preamble
\newtcolorbox[
use counter=shellcounter,
number format=\Alph]{shell}[2][]{%
	colback=blue!5!white,
	colframe=blue!75!white,
	fonttitle=\bfseries,
	title= {\sc Shell \thetcbcounter}: #2,#1}


\newcounter{outputcounter}
\newtcolorbox[
use counter=outputcounter,
number format=\Roman]{result}[2][]{%
	colback=green!5!white,
	colframe=green!75!black,
	fonttitle=\bfseries,
	title= {\sc Result box \thetcbcounter}: #2,#1}


\begin{document}
\begin{titlepage}
\center % Center everything on the page

%-------------------------------------------------------------------------------------
%	HEADING SECTIONS
%------------------------------------------------------------------------------------
\textbf{\large Department of Electronic and Telecommunication Engineering}\\[0.5cm]
\textbf{\Large University of Moratuwa, Sri Lanka}\\[1cm]
\textbf{\large EN4603 - Digital IC Design}\\[2cm]
\includegraphics[width=0.3\textwidth]{figures/uomlogo}\\[2cm]

	
%-------------------------------------------------------------------------------------
%	TITLE SECTION
%------------------------------------------------------------------------------------
\textbf{\Huge Laboratory Experiment 1 \\RTL Synthesis }\\[0.2cm]
\textbf{\Large Laboratory Report}\\[3cm]


%----------------------------------------------------------------------------------------
%	MEMBERS SECTION
%----------------------------------------------------------------------------------------


\vfill

\textbf{\large Submitted by}\\[0.5cm]

\textbf{\large Submitted by}\\[2mm]
\begin{minipage}{0.3\textwidth}
	\begin{flushleft}
		{\large C.S.Pallikkonda}\\[2mm]
		{\large R.M.A.S.Rathnayake }\\[2mm]
		{\large B.P.Thalagala }\\[2mm]		
		
	\end{flushleft}
\end{minipage}
\hspace{2mm}
\begin{minipage}{0.2\textwidth}
	\begin{flushright}
		{\large 180441C }\\[2mm]
		{\large 180534N }\\[2mm]
		{\large 180631J }\\[2mm]

	\end{flushright}
\end{minipage}\\[1cm]

	
	
	
%----------------------------------------------------------------------------------------
%	DATE SECTION
%----------------------------------------------------------------------------------------

\textbf{\large Submitted on}\\[0.5cm]
\textbf{\Large \today} % Date, change the \today to a set date if you want to be precise

%----------------------------------------------------------------------------------------

\vfill % Fill the rest of the page with whitespace

\end{titlepage}

\pagebreak

\tableofcontents
\vfill
\textit{\textbf{Note:}}\\
\textit{All the materials related to the report can also be found at \url{https://github.com/bimalka98/Digital-IC-Design}}

\pagebreak
\listoffigures
\listoftables
\section*{List of Abbreviations}
\addcontentsline{toc}{section}{List of Abbreviations}
\begin{acronym}
	\acro{asic}[ASIC]{Application Specific Integrated Circuit}
	\acro{atpg}[ATPG]{Automatic Test Pattern Generation}
	
	\acro{def}[DEF]{Design Exchange Format}
	\acro{dft}[DFT]{Design For Testability}
		
	\acro{gpdk}[GPDK]{Generic Process Design Kit}
		
	\acro{hdl}[HDL]{Hardware Description Language}
	
	\acro{ic}[IC]{Integrated Circuit}
		
	\acro{rtl}[RTL]{Register-Transfer Level}
	\acro{rx}[RX]{Receiver}
		
	\acro{tcl}[TCL]{Tool Command Language}
	\acro{tx}[TX]{Transmitter}
	
	\acro{uart}[UART]{Universal Asynchronous Receiver Transmitter}
	
\end{acronym}






\pagebreak
\section{Introduction}

\subsection{Practical}
In this practical, you will be using \textit{Cadence Innovus} to place and route the design you synthesized in Laboratory Experiment 2. As inputs to Innovus, you will provide,

\begin{enumerate}
	\item Source Verilog files
	\item Technology libraries provided by the fabrication plant (here, $45~nm$ educational \ac{gpdk} given by Cadence) : {\tt (.lib, .lef, .tch)}
	
	\begin{itemize}
		\item Library Timing (.lib) files specify timing (cell delay, cell transition time, setup and hold time requirement) and power characteristics of standard cells. Slow and fast libraries 	characterize standard cells with maximum and minimum signal delays, which could occur from process variations.

		\item Tch files are binary files that accurately characterize library elements, that include  capacitance and resistance.
		
		\item Library Exchange Format (LEF) specify design rules, metal capacitances, layer information…etc.
	\end{itemize}
	
	\item Multi Mode Multi Corner file (.view)
	\item Constraints file (.sdc)
	\item Scan \ac{def} file (.scandef)
\end{enumerate}

and will obtain the GDSII file as output, which is an industry standard format for exchanging \ac{ic} layout data. You will then analyze, compare, and comment on the placement, cell count, congestion etc. of the design at various stages of the place and route design flow.

\pagebreak
\vfill
\hrule
{\small
\bibliographystyle{IEEEtran}
\bibliography{refer}
}
%---------------------------------------------------------------------------
\end{document}
-
